\documentclass[11pt, letterpaper]{article}

\usepackage[margin=1in]{geometry}
\usepackage{hyperref}
\usepackage{xcolor}
\usepackage{listings}
\usepackage{enumitem}
\usepackage{booktabs}
\usepackage{graphicx}
\usepackage{fancyhdr}
\usepackage{titlesec}
\usepackage{mdframed}
\usepackage{parskip}

% Colors
\definecolor{codeblue}{RGB}{30, 90, 160}
\definecolor{codegray}{RGB}{245, 245, 245}
\definecolor{tipgreen}{RGB}{0, 120, 60}
\definecolor{warnred}{RGB}{180, 30, 30}
\definecolor{noteblue}{RGB}{20, 80, 150}
\definecolor{headergray}{RGB}{60, 60, 60}

% Section formatting
\titleformat{\section}{\large\bfseries\color{headergray}}{}{0em}{}[\titlerule]
\titleformat{\subsection}{\normalsize\bfseries\color{codeblue}}{}{0em}{}

% Code listing style
\lstset{
  backgroundcolor=\color{codegray},
  basicstyle=\ttfamily\small,
  breaklines=true,
  frame=single,
  framerule=0pt,
  rulecolor=\color{codegray},
  xleftmargin=8pt,
  xrightmargin=8pt,
  aboveskip=6pt,
  belowskip=6pt,
}

% Tip box
\newmdenv[
  backgroundcolor=green!8,
  linecolor=tipgreen,
  linewidth=1.5pt,
  leftmargin=0pt,
  rightmargin=0pt,
  innerleftmargin=10pt,
  innerrightmargin=10pt,
  innertopmargin=6pt,
  innerbottommargin=6pt,
  skipabove=8pt,
  skipbelow=8pt,
]{tipbox}

% Note box
\newmdenv[
  backgroundcolor=blue!5,
  linecolor=noteblue,
  linewidth=1.5pt,
  leftmargin=0pt,
  rightmargin=0pt,
  innerleftmargin=10pt,
  innerrightmargin=10pt,
  innertopmargin=6pt,
  innerbottommargin=6pt,
  skipabove=8pt,
  skipbelow=8pt,
]{notebox}

% Warning box
\newmdenv[
  backgroundcolor=red!5,
  linecolor=warnred,
  linewidth=1.5pt,
  leftmargin=0pt,
  rightmargin=0pt,
  innerleftmargin=10pt,
  innerrightmargin=10pt,
  innertopmargin=6pt,
  innerbottommargin=6pt,
  skipabove=8pt,
  skipbelow=8pt,
]{warnbox}

% Header/footer
\pagestyle{fancy}
\fancyhf{}
\fancyhead[L]{\small\textcolor{gray}{Wax Bite Image Processor}}
\fancyhead[R]{\small\textcolor{gray}{User Guide}}
\fancyfoot[C]{\small\thepage}
\renewcommand{\headrulewidth}{0.4pt}

\hypersetup{
  colorlinks=true,
  linkcolor=codeblue,
  urlcolor=codeblue,
}

% -----------------------------------------------------------------------
\begin{document}

\begin{center}
  {\LARGE\bfseries Wax Bite Image Processor}\\[6pt]
  {\large Setup \& User Guide}\\[4pt]
  {\small\textcolor{gray}{macOS and Windows}}
\end{center}

\vspace{0.5em}
\hrule
\vspace{1em}

\tableofcontents

\newpage

% -----------------------------------------------------------------------
\section{Overview}

The Wax Bite Image Processor analyzes wax bite print images of gecko teeth taken at multiple
time points to track tooth eruption and loss over time. Images are processed through a pipeline
of steps, each building on the output of the previous one.

\subsection*{Pipeline Summary}

\begin{center}
\begin{tabular}{lll}
\toprule
\textbf{Step} & \textbf{Command} & \textbf{Description} \\
\midrule
Template Matching & \texttt{match}   & Automatically locates teeth in each image \\
Manual Annotation & \texttt{manual}  & Review and correct automated results \\
Fit \& Project    & \texttt{fitproj} & Fits a hyperbola to the dental arch; projects to 1D \\
Format            & \texttt{format}  & Combines data across time points; generates plots \\
Analyze           & \texttt{analyze} & Opens plots in an interactive exploration window \\
\bottomrule
\end{tabular}
\end{center}

\begin{notebox}
\textbf{Which steps are required?} \texttt{manual} and \texttt{fitproj} are the core required steps.
\texttt{match} is optional but saves time by pre-populating annotations. \texttt{format} and
\texttt{analyze} are run after all images are processed to view results.
\end{notebox}

% -----------------------------------------------------------------------
\section{Folder Structure}

The application is designed to work with one animal's images at a time. A recommended folder
structure keeps each animal's data organized outside the repository so the repository stays clean
between sessions.

\subsection*{Recommended Layout}

\begin{lstlisting}
Documents/
  GeckoData/                        <- your storage folder (name it anything)
    Animals/
      animal1/
        images/                     <- original images live here permanently
        processed/                  <- results moved here after each session
      animal2/
        images/
        processed/
  Wax-Bite-Image-Processor/         <- the repository (cloned once)
    TemplateMatching/
      img/                          <- working folder: copy images here to process
      processed/                    <- working folder: results appear here
      template/                     <- tooth templates for automated matching
      ...
\end{lstlisting}

\subsection*{Workflow for Each Animal}

Follow these steps each time you want to process a new animal's images. Your originals are
never deleted --- you only ever work with copies inside the repository.

\begin{enumerate}[leftmargin=*, label=\textbf{\arabic*.}]
  \item \textbf{Copy images into the repository.}
    Copy all image files from \texttt{Animals/animal1/images/} into
    \texttt{Wax-Bite-Image-Processor/TemplateMatching/img/}.
    Your originals remain safely in the \texttt{Animals} folder.

  \item \textbf{Restore previous results (if re-analyzing).}
    If you have previously processed this animal and want to continue from where you left
    off, copy the contents of \texttt{Animals/animal1/processed/} back into
    \texttt{TemplateMatching/processed/}.

  \item \textbf{Run the analysis.}
    Follow the instructions in Sections~\ref{sec:installation} through~\ref{sec:running}.

  \item \textbf{Move results to the animal's storage folder.}
    When done, move (not copy) the contents of \texttt{TemplateMatching/processed/} into
    \texttt{Animals/animal1/processed/}.

  \item \textbf{Delete the image copies from the repository.}
    Delete all files from \texttt{TemplateMatching/img/}. Your originals are untouched.
    The repository is now clean and ready for the next animal.
\end{enumerate}

\subsection*{Image Naming}

Images must be named in the format \texttt{MM\_DD\_YYYY} followed by any additional
identifier (e.g.\ \texttt{03\_15\_2024\_LG282.jpg}). The date prefix is used to sort
images chronologically and to label results in plots. Images without this format will
cause errors.

% -----------------------------------------------------------------------
\section{One-Time Installation}
\label{sec:installation}

\subsection*{macOS}

\textbf{Step 1: Install Homebrew (if not already installed)}

Homebrew is a package manager that makes installing developer tools straightforward.
Open the Terminal application (find it in Applications $\to$ Utilities, or search with
Spotlight) and paste:

\begin{lstlisting}
/bin/bash -c "$(curl -fsSL https://raw.githubusercontent.com/Homebrew/install/HEAD/install.sh)"
\end{lstlisting}

Follow the on-screen prompts. When it finishes, close and reopen Terminal.

\textbf{Step 2: Install Python and Git}

\begin{lstlisting}
brew install python@3.11 git
\end{lstlisting}

\textbf{Step 3: Clone the repository}

Navigate to the folder where you want to store the repository (e.g.\ \texttt{\textasciitilde/Documents/GeckoData}),
then clone:

\begin{lstlisting}
cd ~/Documents/GeckoData
git clone https://github.com/ecytryn/Wax-Bite-Image-Processor.git
cd Wax-Bite-Image-Processor/TemplateMatching
\end{lstlisting}

\textbf{Step 4: Create and activate a virtual environment}

\begin{lstlisting}
python3.11 -m venv venv
source venv/bin/activate
\end{lstlisting}

You will see \texttt{(venv)} appear at the start of your terminal prompt when the environment
is active.

\textbf{Step 5: Install required packages}

\begin{lstlisting}
pip install -r requirements.txt
\end{lstlisting}

This may take several minutes. When it finishes and the prompt returns, installation is complete.

\textbf{Step 6: Initialize the folder structure}

\begin{lstlisting}
python main.py
\end{lstlisting}

This creates the \texttt{img/}, \texttt{template/}, and \texttt{processed/} subfolders.

\begin{tipbox}
\textbf{Verifying the installation:} After running \texttt{python main.py} for the first time
you should see an \texttt{img/} folder appear inside \texttt{TemplateMatching/}. If you see
error messages instead, check the Troubleshooting section.
\end{tipbox}

% ---------------------------------------------------
\subsection*{Windows}

\textbf{Step 1: Install Python}

\begin{enumerate}
  \item Go to \url{https://www.python.org/downloads/release/python-3110/} in your browser.
  \item Scroll down and click \textbf{Windows installer (64-bit)} to download.
  \item Run the installer. \textbf{Important:} before clicking ``Install Now'', check the
        box that says \textbf{``Add Python to PATH''} at the bottom of the window. This
        step is critical. If you have administrator privileges, also check
        \textbf{``Install for all users''} so that everyone who connects to this machine
        can use the app.
  \item Click ``Install Now'' and wait for it to finish.
\end{enumerate}

\begin{tipbox}
To verify Python installed correctly: press the Windows key, type \texttt{cmd}, press
Enter. In the black window that opens, type \texttt{python --version} and press Enter.
You should see \texttt{Python 3.11.x}.
\end{tipbox}

\textbf{Step 2: Install Git}

\begin{enumerate}
  \item Go to \url{https://git-scm.com/download/win} --- the download starts automatically.
  \item Run the installer, clicking Next through all steps; the defaults are fine.
\end{enumerate}

\textbf{Step 3: Clone the repository}

\begin{enumerate}
  \item Open File Explorer and navigate to the folder where you want to store the
        repository (e.g.\ your Documents folder).
  \item Click in the address bar at the top of File Explorer, type \texttt{cmd}, and
        press Enter. A Command Prompt window will open in that folder.
  \item Type the following and press Enter:
\end{enumerate}

\begin{lstlisting}
git clone https://github.com/ecytryn/Wax-Bite-Image-Processor.git
cd Wax-Bite-Image-Processor\TemplateMatching
\end{lstlisting}

\textbf{Step 4: Create and activate a virtual environment}

\begin{lstlisting}
python -m venv venv
venv\Scripts\activate.bat
\end{lstlisting}

You will see \texttt{(venv)} appear at the start of the line when the environment is active.

\textbf{Step 5: Install required packages}

\begin{lstlisting}
pip install -r requirements.txt
\end{lstlisting}

\textbf{Step 6: Initialize the folder structure}

\begin{lstlisting}
python main.py
\end{lstlisting}

\begin{warnbox}
\textbf{Every session:} You must activate the virtual environment each time you open a new
terminal window before running the application. On macOS: \texttt{source venv/bin/activate}.
On Windows: \texttt{venv\textbackslash Scripts\textbackslash activate.bat}. If you forget,
you will see a \texttt{ModuleNotFoundError}.
\end{warnbox}

% -----------------------------------------------------------------------
\section{Running the Application}
\label{sec:running}

All commands are run from inside the \texttt{TemplateMatching/} folder with the virtual
environment active.

\subsection*{Basic Command Format}

\begin{lstlisting}
python main.py [steps] [image selection flags] [display flags]
\end{lstlisting}

\textbf{Steps} (one or more, space-separated): \texttt{match}, \texttt{manual},
\texttt{fitproj}, \texttt{format}, \texttt{analyze}

\textbf{Image selection flags:}
\begin{center}
\begin{tabular}{ll}
\toprule
\textbf{Flag} & \textbf{Meaning} \\
\midrule
\texttt{-s N} & Start at image index \texttt{N} (first image is index 0) \\
\texttt{-n N} & Process \texttt{N} images \\
\bottomrule
\end{tabular}
\end{center}

\textbf{Display flags:}
\begin{center}
\begin{tabular}{ll}
\toprule
\textbf{Flag} & \textbf{Meaning} \\
\midrule
\texttt{--width N} & Set annotation window width to \texttt{N} pixels for this session \\
\bottomrule
\end{tabular}
\end{center}

The default window width is controlled by \texttt{MAX\_WIDTH} in \texttt{utils.py} (default: 800 pixels).
Use \texttt{--width} to override this for a session without editing the file. If the image appears
too small or too large, try adjusting this value. A good starting point is roughly half your screen
width. The window position on screen can be adjusted by editing \texttt{WINDOW\_X} and \texttt{WINDOW\_Y}
in \texttt{utils.py}.

\subsection*{Common Commands}

\begin{lstlisting}
# Run the full pipeline on all images
python main.py

# Run only template matching on all images
python main.py match

# Run manual annotation on all images
python main.py manual

# Run manual annotation with a custom window width
python main.py manual --width 1200

# Run fit and project on all images
python main.py fitproj

# Run multiple steps in sequence on all images
python main.py match manual fitproj

# Process only the 3rd image (index 2) through manual annotation
python main.py -s 2 -n 1 manual

# Process 5 images starting from the 1st
python main.py -s 0 -n 5 match fitproj

# Generate and view plots (no image selection needed)
python main.py format
python main.py analyze

# Process specific named images
python main.py 03_15_2024_LG282.jpg 03_22_2024_LG282.jpg manual
\end{lstlisting}

\begin{notebox}
\textbf{Recommended workflow for a new dataset:}
\begin{enumerate}[leftmargin=*, topsep=2pt, itemsep=2pt]
  \item \texttt{python main.py match} --- run automated matching on all images
  \item \texttt{python main.py manual} --- manually review and correct each image
  \item \texttt{python main.py fitproj} --- fit hyperbola and project
  \item \texttt{python main.py analyze} --- view results interactively
\end{enumerate}
Running steps separately (rather than all at once) gives you more control and makes it
easier to identify which step caused an error.
\end{notebox}

% -----------------------------------------------------------------------
\section{Manual Annotation (the \texttt{manual} step)}

When you run the \texttt{manual} step, a window opens showing the current image. If
template matching has already been run, those results are pre-loaded as a starting point.

\subsection*{Annotation Modes}

There are four types of markers you can place:

\begin{center}
\begin{tabular}{lll}
\toprule
\textbf{Mode} & \textbf{Key} & \textbf{Use for} \\
\midrule
Tooth        & \texttt{1} & A tooth \\
Gap          & \texttt{2} & A gap between teeth \\
Center Tooth & \texttt{3} & The central-most tooth (only one allowed) \\
Center Gap   & \texttt{4} & The central-most gap (only one allowed) \\
\bottomrule
\end{tabular}
\end{center}

\subsection*{Keyboard Shortcuts}

\begin{center}
\begin{tabular}{ll}
\toprule
\textbf{Key} & \textbf{Action} \\
\midrule
\texttt{1} / \texttt{2} / \texttt{3} / \texttt{4} & Switch to tooth / gap / center tooth / center gap mode \\
\texttt{Tab}        & Cycle through modes in order \\
\texttt{Space}      & Toggle boxes on/off (view image without overlays) \\
\texttt{S}          & Save and close current image \\
\texttt{Q}          & Quit without saving \\
\texttt{Esc}        & Close current image without saving \\
$\leftarrow$ arrow  & Save and go to previous image \\
$\rightarrow$ arrow & Save and go to next image \\
\bottomrule
\end{tabular}
\end{center}

\subsection*{Mouse Controls}

\begin{itemize}[leftmargin=*]
  \item \textbf{Left click} on empty space: place a marker of the current mode
  \item \textbf{Left click} inside an existing box: delete that marker
\end{itemize}

\begin{warnbox}
\textbf{Always save before closing} using \texttt{S} or the arrow keys. \texttt{Q} and
\texttt{Esc} exit without saving the \emph{current} image, which will cause \texttt{fitproj}
to fail for that image if no previous annotation exists for it. Images you navigated
through using the arrow keys are saved automatically and are not affected.
\end{warnbox}

\begin{tipbox}
\textbf{Window size:} The annotation window opens at the width set by \texttt{MAX\_WIDTH}
in \texttt{utils.py} (default: 800 pixels), or the value passed via \texttt{--width} on
the command line. If the image appears too small, try \texttt{python main.py manual --width 1200}
or a larger value. Avoid resizing the window manually by dragging its edges, as this will
cause click positions to be misaligned with markers.
\end{tipbox}

% -----------------------------------------------------------------------
\section{Tooth Templates}

The \texttt{template/} folder contains small image patches of individual teeth used
by the \texttt{match} step to automatically locate teeth in each image. The current
templates were created from images produced with the standard staining protocol used
in the lab.

\subsection*{When to add templates}

If you notice that a consistent type or shape of tooth is frequently missed by the
automated matching, adding one or more example images of that tooth to the
\texttt{template/} folder will help the matcher find it. Templates should be:

\begin{itemize}[leftmargin=*]
  \item Small crops of a single tooth (not multiple teeth)
  \item Representative of the appearance you want to detect
  \item Saved as \texttt{.jpg} or \texttt{.png}
\end{itemize}

\subsection*{Using unstained images}

If you are working with unstained wax bites, the existing templates are unlikely to
produce good results since tooth appearance will differ significantly. In this case,
create a new set of templates from your unstained images and either replace or
supplement the existing ones in the \texttt{template/} folder. The \texttt{manual}
step can always be used to correct or replace the automated results regardless of
template quality.

% -----------------------------------------------------------------------
\section{Interactive Analysis (the \texttt{analyze} step)}

Running \texttt{python main.py analyze} generates four plots and opens them in
interactive windows:

\begin{itemize}[leftmargin=*]
  \item \textbf{Arclength plot} --- tooth positions in arclength coordinates over time
  \item \textbf{Index plot} --- tooth positions by index over time
  \item \textbf{Eruption plot} --- eruption events colored by time since last eruption
  \item \textbf{Fall-out plot} --- tooth loss events colored by time since last loss
\end{itemize}

\subsection*{Interacting with the plots}

\begin{itemize}[leftmargin=*]
  \item \textbf{Left click} on a point: opens the projected 1D image for that date
  \item \textbf{Left click and drag} across a range of dates: opens the projected 1D
        images for all dates in the selected range, stacked vertically and aligned to
        the same width. This is useful for visually comparing annotations across nearby
        days and identifying errors such as a tooth misidentified as a gap, or a missed
        gap causing all teeth distal to it to appear shifted relative to other days.
  \item \textbf{Right click} (ctrl+click on Mac): opens the manual annotation tool for
        that date. Note: changes will not appear in the plots until you re-run
        \texttt{format} or \texttt{analyze}.
\end{itemize}

\begin{tipbox}
The matplotlib plot windows also have a built-in toolbar that allows you to zoom and
pan within the plots. Look for the toolbar at the bottom of each plot window.
\end{tipbox}

% -----------------------------------------------------------------------
\section{Troubleshooting}

\begin{description}[leftmargin=0pt, style=nextline]

  \item[\texttt{ModuleNotFoundError} when running the app]
    The virtual environment is not activated. Make sure you see \texttt{(venv)} at
    the start of your command line.
    macOS: \texttt{source venv/bin/activate} \quad
    Windows: \texttt{venv\textbackslash Scripts\textbackslash activate.bat}

  \item[\texttt{AssertionError: manual data.csv does not exist}]
    The \texttt{fitproj} step cannot find annotation data for one or more images.
    This usually means the \texttt{manual} step was exited with \texttt{Esc} or
    \texttt{Q} rather than saved with \texttt{S} or the arrow keys. Re-run
    \texttt{python main.py manual} and save each image with \texttt{S}.

  \item[\texttt{python} not recognized (Windows)]
    Python was not added to PATH during installation. Uninstall Python and reinstall,
    making sure to check ``Add Python to PATH'' at the bottom of the installer window.

  \item[Image appears too small or too large]
    Use the \texttt{--width} flag to set a comfortable window size, for example:
    \texttt{python main.py manual --width 1200}. A good starting point is roughly half
    your screen width. Once you find a value that works, update \texttt{MAX\_WIDTH} in
    \texttt{utils.py} to make it the default.

  \item[Click positions are misaligned with markers]
    This can happen if the annotation window has been manually resized by dragging its
    edges. Close the window and reopen it without resizing. If the problem persists,
    try a different \texttt{--width} value. This issue is particularly common over
    remote desktop connections.

  \item[Python not found when connecting remotely (Windows)]
    Python may have been installed for a specific user account rather than for all users.
    Ask an administrator to reinstall Python with the ``Install for all users'' option
    checked.

  \item[\texttt{RuntimeError: Unable to fit a Hyperbola or Parabola}]
    The \texttt{fitproj} step detected a circle or ellipse instead of a hyperbola.
    This is usually caused by annotation errors. Open the image in the \texttt{manual}
    step, check for misplaced markers (especially any that are far from the dental arch),
    correct them, save, and re-run \texttt{fitproj}.

  \item[A window appears and immediately closes]
    Look at the terminal for an error message. If you see a traceback (lines beginning
    with \texttt{File "..."}), note the last few lines --- they identify the specific
    error. Refer to this guide or contact the developer.

  \item[Plots are empty or show no data]
    Make sure you have run \texttt{fitproj} on all images before running \texttt{format}
    or \texttt{analyze}. Also check that image filenames follow the \texttt{MM\_DD\_YYYY}
    format.

\end{description}

% -----------------------------------------------------------------------
\section*{Quick Reference}

\begin{center}
\begin{tabular}{ll}
\toprule
\textbf{Command} & \textbf{What it does} \\
\midrule
\texttt{python main.py} & Run full pipeline on all images \\
\texttt{python main.py match} & Template matching, all images \\
\texttt{python main.py manual} & Manual annotation, all images \\
\texttt{python main.py manual --width 1200} & Manual annotation with custom window width \\
\texttt{python main.py fitproj} & Fit \& project, all images \\
\texttt{python main.py format} & Generate plots \\
\texttt{python main.py analyze} & Interactive plot viewer \\
\texttt{python main.py -s 2 -n 1 manual} & Manual annotation, image index 2 only \\
\texttt{python main.py match fitproj} & Match then fitproj, all images \\
\bottomrule
\end{tabular}
\end{center}

\vspace{1em}
\begin{center}
\small\textcolor{gray}{Source code: \url{https://github.com/ecytryn/Wax-Bite-Image-Processor}}
\end{center}

\end{document}
